\chapter{平方元数：一套可以写成公式的控制器}\label{ch:square}
\section{为什么尝试平方}
令 $k=t^2$，总位置数取 $n=t^2+t+1$，其中 $t\ge2$。核心形状是
\[K(u,v)=(u^{t^2},v^t,w),\qquad w=-t^2u-tv.
\]
平方给出一个自然的库存比例：大块恰好能分成 $t$ 组，每组 $t$ 个位置。后面会看到，这个比例还让某个返回矩阵满足三次周期。

这一讲先解决已经进入核心以后的控制问题，下一讲再解决任意输入的入口和最后的算术条件。这样可以清楚看到：哪些是矩阵恒等式，哪些还需要关于整数的定理。\src{S13}

始终取差坐标 $x=u,y=u-v$。本原化参数后，$G(K)=\gcd(y,n)$，而 $\gcd(n,t)=1$，所以合法核心恰好满足 $\gcd(y,n)=1$。

\section{三个一步返回}
分别平均
\[(u^{t^2-t},v^t),\qquad(u^{t^2-1},w),\qquad
(u^{t^2-t},v^{t-1},w).
\]
第一种留下 $t$ 份旧 $u$，第二种留下 $t$ 份旧 $v$，第三种留下 $t$ 份旧 $u$；其余单点也明确保留。三个实际作用为
\[
\mathcal A=\frac1t A,\quad A=\begin{pmatrix}t&-1\\0&-1\end{pmatrix},\qquad
\mathcal B=\frac1{t^2}B,\quad B=\begin{pmatrix}-t-1&t\\-n&t^2+t\end{pmatrix},
\]
\[
\mathcal C=\frac1{t^2}C,\qquad C=\begin{pmatrix}-t-1&1\\-n&1\end{pmatrix}.
\]
例如第三种的新大块均值为 $-(tu+v)/t^2$，新小块值为 $u$。相减后便得到 $C$ 的第二行。这说明矩阵并非猜出来的：每个系数都能从选取的位置数直接读出。

\section{周期把不可逆平均变成可逆的方向控制}
矩阵 $C$ 的迹是 $-t$、行列式是 $t^2$，特征恒等式给出
\[C^2+tC+t^2I=0,\qquad C^3=t^3I.
\]
于是实际返回满足 $\mathcal C^3=t^{-3}I$。连续三次返回缩小了状态，却把方向完全还原。这正是可以实现射影逆的原因。

另外，$\operatorname{tr}(CA)=0$，从而
\[(\mathcal C\mathcal A)^2=t^{-3}I.
\]
因此 $\mathcal A$ 也有正向射影逆。最后
\[C^2B=t^3\Sigma,\qquad
\Sigma=\begin{pmatrix}1&-1\\0&-1\end{pmatrix}.
\]
至此，$\mathcal A,\mathcal B,\mathcal C,\Sigma$ 的射影逆都有真正的向前执行方式。

因为 $\Sigma\mathcal A=\Delta(1/t)$，我们已有 $\Delta(t)^{\pm1}$。交换子给出
\[\Sigma\Delta(t)\Sigma\Delta(t)^{-1}=U((t-1)/t).
\]
取幂和共轭后，得到 $U((t-1)R_0)$，其中 $R_0=\Z[1/t]$。借助 $C$ 的横向作用，前篇的同余子群工具给出一个深层
\[\Gamma((t-1)^2n^3,R_0).
\]
这里先出现的是带因子 $t-1$ 的剪切。这足以准备高精度，却还没有得到任意大小的上剪切。

\section{两步返回供应所有小整数倍率}
选 $0\le j\le t-1$，先平均
\[(u^{t^2-j-1},v^j,w),
\]
得到 $t^2$ 份新值。保留其中 $t$ 份，以及一个未选的旧 $u$，再平均其余 $t^2$ 个位置。实际矩阵是 $M_j/t^3$，其中
\[M_j=\begin{pmatrix}1&j-t\\n&(t+1)(j-t)\end{pmatrix}.
\]
所有选择都在库存内：第一次之后至少剩一份旧 $u$，新值则足够保留 $t$ 份。

最有用的地方是这些矩阵第一列相同。因而
\[M_0^{-1}M_j=\Delta((t-j)/t).
\]
$\det M_0=t^3$，是 $R_0$ 中的单位；它在相应有限商中保持已取得的深层。应用前篇的有限商幂论证，可以用正向词实现 $M_0$ 的射影逆。配合 $\Delta(t)$，便得到
\[\Delta(d),\qquad 1\le d\le t.
\]
这一步首先给正向倍率。某个 $d$ 若含有 $n$ 的素因子，其逆会破坏合法域，所以不能一概立即使用。

\section{为什么必须专门处理除以二}
我们希望把旧剪切中的 $t-1$ 除掉。单纯把更多素数加入环并不总能做到，因为 $t-1$ 可能还含不可逆因子。有效办法是重新计算一个交换子：一旦 $\Delta(2)$ 和逆都可执行，就有
\[\Sigma\Delta(2)\Sigma\Delta(1/2)=U(1/2).
\]
新的系数已经不带 $t-1$。这是新恒等式提供的改进。

若 $t$ 为偶数，二已经是 $R_0$ 的单位，初始深层的水平也是奇数，可以用有限商恢复倍率的逆。若 $t$ 为奇数，二可能碰到旧水平中的 $t-1$，要用一个专门的行列式一矩阵：
\[
H=\frac1{2t}\Delta(2)C\Delta(2)
=\begin{pmatrix}
-(t+1)/(2t)&1/t\\-n/t&2/t
\end{pmatrix}\in\SL_2(R_0).
\]
它属于 $\SL_2(R_0)$ 的关键是 $t+1$ 为偶数，且行列式等于一。$H$ 本身由已实现因子射影给出；它的某个正整数次幂落入深层，所以有正向射影逆。再用
\[C\Delta(2)H^{-1}=2t\Delta(1/2)
\]
就实现了所需的除二倍率。

于是，令 $R=\Z[1/(2t)]$，我们得到 $U(R)$。下剪切来自另一个直接恒等式：
\[U(-1)\mathcal C\Delta(t^2)=L(-n/t^2).
\]
取幂、共轭，得到全部 $L(nR)$。平方控制器的主要工作至此完成：原位置平均已经供应了一整套上剪切、$n$ 倍下剪切和若干倍率。

\section{继续扩大允许倍率的范围}
对于 $1\le d\le t$ 且 $\gcd(d,n)=1$ 的整数，$\Delta(d)$ 已有正向实现。现在用新的同余层，可以恢复它的逆，进而把 $d$ 的素因子加入局部化环。记所得环为 $R_t$。

这里扩大的是\emph{核心方向矩阵}的系数环。真实坐标仍只有 $k$ 所允许的分母。例如 $\Delta(1/5)$ 与 $\diag(5,1)$ 表示同一个射影作用；实际返回还带一个由平均词决定的共同尺度。因而方向矩阵出现五的分母，不表示许可操作突然变成了五平均。这正是前篇共同尺度约定在此处的用途。

可逆倍率越来越多，模 $n$ 能改变的余数也越来越多。不过“很多”还不等于“全部”。下一讲会把唯一剩下的乘法问题写成一个有限群等式。至于本讲的周期、除二和剪切公式，都对每个 $t\ge2$ 成立，不依赖 GRH。
